\documentclass[10pt,conference]{IEEEtran}
\usepackage{cite}
\usepackage{amsmath,amssymb,amsfonts}
\usepackage{algorithmic}
\usepackage{graphicx}
\usepackage{textcomp}
\usepackage{xcolor}
\usepackage{booktabs}
\usepackage{url}
\usepackage{hyperref}

\hypersetup{
    colorlinks=true,
    linkcolor=blue,
    citecolor=blue,
    urlcolor=blue
}

\def\BibTeX{{\rm B\kern-.05em{\sc i\kern-.025em b}\kern-.08em
    T\kern-.1667em\lower.7ex\hbox{E}\kern-.125emX}}

\begin{document}

\title{RecoveryEval: Evaluating Continuation Inertia Under Injected Reasoning Errors in Language Models}

\author{\IEEEauthorblockN{Sham Satish Thakare}
\IEEEauthorblockA{\textit{Independent Researcher} \\
Pune, Maharashtra, India \\
Email: shamthakare@gmail.com}
}

\maketitle

\begin{abstract}
Evaluating whether post-training enhances a language model's capacity to detect and recover from intermediate reasoning errors is critical for building reliable artificial intelligence systems. However, naive trajectory evaluations routinely conflate error recovery with baseline fluency, prompt length discrepancies, and truncation censoring. In this paper, we present \textbf{RecoveryEval}, a state-matched, provenance-aware evaluation framework that assesses reasoning trajectories under paired counterfactual prefix perturbations. We evaluate 400 genuine neural rollouts generated by two released checkpoints (\texttt{Qwen2.5-Math-1.5B} Base and Instruct) across 20 multi-step arithmetic problems under a length-balanced, truncation-free protocol (Protocol 2.0; 256-token budget). 

Under symmetrically matched prefixes withholding the terminal answer, neither model exhibits spontaneous error recovery: across all 200 recovery continuations, zero spontaneous corrections are observed ($0/200$, $0.0\%$). Instead, models exhibit pronounced \textbf{continuation inertia}: in 98.5\% of cases (197/200), models accept the injected error as an authoritative premise and continue derivations consistently from it, with the remaining 1.5\% failing due to attention slips. Under an independent-rollout assumption ($n=200$), the exact two-sided 95\% Clopper-Pearson confidence interval is $[0.00\%, 1.83\%]$ (one-sided upper bound: $1.49\%$); accounting for intra-cluster correlation across the 20 problem templates yields an effective cluster-level sensitivity bound of $\le 13.91\%$ (two-sided interval: $[0.00\%, 16.84\%]$). On unperturbed control prefixes, both models achieve 100.0\% structural accuracy, yielding an observed sample recovery contrast of $D = 0.0000$ (where sample variance collapses due to boundary saturation across the evaluated items). Trace forensics reveal that legacy regex parsers falsely penalize base models by 3\% due to trailing review summaries, which structural equivalence resolves. Our findings demonstrate that small language models behave as faithful continuation engines rather than autonomous truth-checkers, highlight how naive evaluation designs manufacture severe confounding artifacts, and provide open-source cryptographic infrastructure for reasoning diagnostics.
\end{abstract}

\begin{IEEEkeywords}
Language Models, Mathematical Reasoning, Error Recovery, State-Matched Evaluation, Benchmark Governance, Continuation Inertia, Clopper-Pearson Bounds.
\end{IEEEkeywords}

\section{Introduction}
Large language models (LLMs) specialized for mathematics have demonstrated remarkable gains on multi-step reasoning benchmarks \cite{cobbe2021gsm8k, qwen25math2024}. A central question in AI reliability is whether post-training procedures (such as instruction tuning, reinforcement learning from human feedback, or process supervision) instill the capacity to recognize and self-correct intermediate calculation mistakes \cite{lightman2023process, zelikman2022star}. 

However, measuring recovery behavior is fraught with methodological pitfalls. Naive comparisons between unperturbed and error-injected trajectories frequently introduce severe confounding variables:
\begin{enumerate}
    \item \textbf{Control Answer Exposure}: If control prefixes contain the ground-truth final answer while recovery prefixes withhold it, the evaluation measures pattern echoing rather than reasoning.
    \item \textbf{Truncation Censoring}: Tight generation limits (e.g., 64 tokens) prematurely truncate longer reasoning traces, confounding verbosity with error recovery.
    \item \textbf{Evaluator Brittleness}: Regular-expression parsers that extract the last numerical token falsely penalize models that conclude with conversational summary steps.
    \item \textbf{Pseudoreplication}: Treating clustered rollouts as independent experimental units inflates degrees of freedom and fabricates statistical significance.
\end{enumerate}

In preliminary experiments using an unbalanced 64-token protocol (Protocol 1.0), an apparent negative interaction was reported ($D \approx -0.21$), suggesting that instruction tuning degraded recovery resilience. However, Protocol 1.0 suffered simultaneously from 100\% control answer leakage, 53.0\% truncation censoring, and decorative matching formulas.

To resolve these ambiguities, we present \textbf{RecoveryEval}, an open-source evaluation framework built on \textit{Paired Counterfactual Prefix Perturbation Evaluation} (PCPPE) and cryptographic event ledgers. Across 400 genuine neural rollouts evaluated under a length-balanced, 256-token budget (Protocol 2.0), we establish:
\begin{enumerate}
    \item \textbf{Non-Reproduction of Negative Interaction}: Under Protocol 2.0, the historical negative interaction disappears ($D = 0.0000$ under structural verification; $D = -0.0300, p = 0.0828$ under heuristic parsing), proving that recovery estimates are acutely sensitive to protocol design.
    \item \textbf{Continuation Inertia}: Small (1.5B) language models do not spontaneously self-correct from injected intermediate errors ($0/200$ recoveries). In 98.5\% of cases, they treat corrupted premises as authoritative and faithfully propagate errors downstream.
    \item \textbf{Finite-Sample Uncertainty Bounds}: We establish that observing zero recoveries corresponds to a one-sided 95\% upper bound of 1.49\% under independent rollout assumptions ($n=200$), but cluster-level modeling reveals an effective sensitivity upper bound of 13.91\% ($C=20$).
    \item \textbf{Evaluator Sensitivity}: Heuristic regex extractors artificially penalize base models by 3\% due to trailing review summaries, whereas structural equivalence demonstrates 100\% baseline accuracy for both checkpoints.
\end{enumerate}

\section{Related Work}
Self-correction and error localization in language models have attracted intense scrutiny. Table \ref{tab:literature} positions \textbf{RecoveryEval} relative to existing paradigms.

\begin{table*}[t]
\caption{Systematic Comparison of Reasoning and Self-Correction Evaluation Frameworks}
\label{tab:literature}
\centering
\begin{tabular}{p{3.2cm}p{3.0cm}p{2.6cm}p{2.8cm}p{4.2cm}}
\toprule
\textbf{Framework} & \textbf{Primary Target} & \textbf{Intervention Type} & \textbf{Prefix Symmetry} & \textbf{Core Empirical Finding} \\
\midrule
GSM8K \cite{cobbe2021gsm8k} & End-to-end reasoning & None (Question only) & None & Scaling improves final outcome accuracy \\
PRM800K \cite{lightman2023process} & Process supervision & Step classification & Unmatched & Step-level oversight outperforms outcome oversight \\
Reflexion \cite{shinn2023reflexion} & Multi-turn correction & Verbal reflection buffer & Asymmetric & External verbal feedback enables error recovery \\
CRITIC \cite{gou2024critic} & Tool-based correction & External interpreter/search & Interactive & Tool feedback enables calculation self-correction \\
Huang et al. \cite{huang2023large} & Intrinsic self-correction & Prompted reflection & Post-hoc & Prompted reflection without feedback degrades accuracy \\
BIG-Bench Mistake \cite{tyen2024bigbenchmistake} & Error localization & Explicit error location cue & Post-hoc traces & Models cannot find errors, but fix given location \\
ProcessBench \cite{zheng2024processbench} & Error step identification & Benchmark critic tasks & Varied & PRMs struggle to locate errors on competition math \\
\midrule
\textbf{RecoveryEval (Ours)} & \textbf{In-flight continuation inertia} & \textbf{Symmetrical step perturbation} & \textbf{Length-matched ($\Delta_{\text{tokens}}=0$)} & \textbf{98.5\% error propagation; 0\% spontaneous recovery} \\
\bottomrule
\end{tabular}
\end{table*}

\subsection{Intrinsic Self-Correction Limitations}
Huang et al. \cite{huang2023large} established that prompting models to self-correct (*"Review your answer and correct mistakes"*) without external feedback impairs performance because models alter correct answers more frequently than they repair incorrect ones. Stechly et al. \cite{stechly2024selfverification} demonstrated similar self-verification failures in planning tasks. While these studies evaluated post-hoc multi-turn prompting, \textbf{RecoveryEval} investigates \textit{in-flight intermediate continuation}: how autoregressive models behave when continuing from corrupted intermediate derivations without explicit reflection prompts.

\subsection{External Feedback and Process Supervision}
Shinn et al. \cite{shinn2023reflexion} (Reflexion) and Gou et al. \cite{gou2024critic} (CRITIC) demonstrated that self-correction succeeds when models interact with external tools (code interpreters, calculators, search engines). Tyen et al. \cite{tyen2024bigbenchmistake} (BIG-Bench Mistake) showed that while models fail to locate their own mistakes, they can correct them if the exact erroneous step is pointed out. Zheng et al. \cite{zheng2024processbench} (ProcessBench) and Wang et al. \cite{wang2024mathshepherd} (Math-Shepherd) confirmed the difficulty of automated process verification. \textbf{RecoveryEval} complements these studies by providing a clean baseline for unassisted continuation, demonstrating that without external tool feedback, continuation inertia completely dominates.

\section{Research Questions and Causal Framing}
We formulate three core empirical questions:
\begin{itemize}
    \item \textbf{RQ1 (Continuation vs. Recovery)}: When presented with an unprompted continuation prefix containing an erroneous intermediate calculation, does an autoregressive language model detect and correct the error, or does it propagate the corrupted premise downstream?
    \item \textbf{RQ2 (Protocol Sensitivity)}: Does the apparent interaction between model alignment and error recovery reproduce when control answer leakage, truncation censoring, and parser artifacts are eliminated?
    \item \textbf{RQ3 (Uncertainty Quantification)}: When zero recovery events are observed in clustered factorial data, what exact finite-sample confidence limits are statistically defensible?
\end{itemize}

\subsection{Causal Identification Boundaries}
We represent the data-generating process as a structural causal model:
\begin{itemize}
    \item \textbf{Prefix Intervention ($T \to Y$)}: Within each problem cluster $P_i$, the problem text and remaining token horizon are held constant. The only intervention is the substitution of the Step 1 product ($s_C$ vs. $s_R$). The contrast in continuation accuracy between $s_C$ and $s_R$ is \textit{causally identified}.
    \item \textbf{Checkpoint Comparison ($M \to Y$)}: Comparing Base and Instruct models is \textit{observational}. Base models expect raw completion strings, while Instruct models require ChatML formatting (\texttt{<|im\_start|>}). We do not claim to isolate a pure causal post-training effect free from interface confounding.
    \item \textbf{Scope of Construct}: The experiment measures \textit{continuation inertia under externally supplied corruption}, not an intrinsic inability to self-correct during uninterrupted, self-generated reasoning.
\end{itemize}

\section{Experimental Design: Protocol 2.0}

\subsection{Benchmark Construction}
We evaluate 20 multi-step arithmetic word problems inspired by GSM8K ($C=20$). Each problem involves an initial rate-multiplication step (Step 1) followed by a downstream multiplication step to determine total revenue (Step 2). Neither prefix reveals the final terminal answer.

\subsection{Symmetrical Counterfactual Prefix Construction}
For each problem $P_i$, two prefixes are constructed:
\begin{itemize}
    \item \textbf{Reference Control Prefix ($s_C$)}: Contains valid Step 1 arithmetic (e.g., $5 \times 4 = 20$).
    \item \textbf{Recovery Perturbation Prefix ($s_R$)}: Introduces an arithmetic corruption into Step 1 (e.g., $5 \times 4 = 22$).
\end{itemize}
Both prefixes are truncated immediately prior to the Step 2 calculation, ensuring exact token and character balance ($\Delta_{\text{tokens}} = 0, \Delta_{\text{chars}} = 0$).

\subsection{Factorial Matrix and Hardware Execution}
The factorial design comprises:
$$20 \text{ problems} \times 2 \text{ models} \times 2 \text{ conditions} \times 5 \text{ seeds} = 400 \text{ rollouts}$$
Evaluated checkpoints are \texttt{Qwen/Qwen2.5-Math-1.5B} (Base, Rev \texttt{4a83ca6e}) and \texttt{Qwen/Qwen2.5-Math-1.5B-Instruct} (Instruct, Rev \texttt{aafeb0fc}) \cite{qwen25math2024}. Rollouts were executed on an NVIDIA GeForce RTX 4090 GPU in native \texttt{float16} precision using stochastic nucleus sampling ($T=0.7, p=0.9$) across 5 fixed seeds ($401, 402, 403, 404, 405$) with a 256-token budget.

\subsection{Cryptographic Provenance and Artifact Integrity}
Every rollout record captures exact input token IDs, generated BPE token IDs, wall-clock inference duration, and model weight digest in an append-only event ledger. The canonical dataset \texttt{10\_COMPLETE\_QWEN\_PROTOCOL2\_ROLLOUTS.jsonl} is frozen under SHA-256 digest \texttt{b8792f8df375fccdbd800d440cc884ef\allowbreak 2db2f860ff3533af01c70161a5cc1d15}.

\section{Statistical Methodology and Results}

\subsection{Cluster-Robust Factorial Outcomes}
Table \ref{tab:outcomes} summarizes completion accuracies across all four experimental cells under heuristic regex and structural semantic verification.

\begin{table}[htbp]
\caption{Protocol 2.0 Factorial Outcomes ($N=400$ Rollouts, $C=20$ Clusters)}
\label{tab:outcomes}
\centering
\begin{tabular}{lccc}
\toprule
\textbf{State Condition} & \textbf{Base ($\pi_0$)} & \textbf{Instruct ($\pi$)} & \textbf{Difference ($\Delta$)} \\
\midrule
Recovery States ($s_R$) & 0.0000 & 0.0000 & $+0.0000$ \\
Control States (Heuristic) & 0.9700 & 1.0000 & $+0.0300$ \\
Control States (Structural) & 1.0000 & 1.0000 & $+0.0000$ \\
\midrule
\textbf{Matched Contrast (Heuristic)} & -- & -- & \textbf{$-0.0300$} \\
\textbf{Cluster SE / $p$-value} & -- & -- & $SE = 0.0164, p = 0.0828$ \\
\textbf{95\% Confidence Interval} & -- & -- & $[-0.0643, +0.0043]$ \\
\midrule
\textbf{Matched Contrast (Structural)} & -- & -- & \textbf{$+0.0000$ (Sample Null)} \\
\textbf{Cluster SE / $p$-value} & -- & -- & Degenerate ($s_D = 0.0000$) \\
\bottomrule
\end{tabular}
\end{table}

\subsection{Degenerate Variance Analysis}
Under structural scoring, every problem cluster $j \in \{1, \dots, 20\}$ yields:
$$\Delta_{j, \text{Base}} = -1.0000, \quad \Delta_{j, \text{Instruct}} = -1.0000 \implies D_j = 0.0000$$
Consequently, the sample cluster standard deviation is $s_D = 0.0000$, and the cluster-robust standard error is $\mathrm{SE} = 0.0000$. Under this condition:
\begin{itemize}
    \item The Student-$t$ statistic $t = \bar{D} / \mathrm{SE} = 0/0$ is \textbf{mathematically undefined}.
    \item The conventional two-sided $p$-value is \textbf{not computable}.
    \item The conventional Studentized 95\% confidence interval $[0, 0]$ is \textbf{degenerate}.
\end{itemize}
We emphasize that zero sample variance across 20 items is an empirical consequence of boundary saturation (100\% control vs. 0\% recovery), rather than proof of universal population equivalence.

\subsection{Verifier Forensics and Heuristic Regex Failures}
In heuristic scoring, Base Control scored 97.0\% due to 3 apparent failures under seed 405 (problems \texttt{fresh\_004}, \texttt{fresh\_006}, \texttt{fresh\_013}). Trace-by-trace inspection revealed that the base model correctly calculated the revenue but concluded with an unprompted summary review ending in list bullets (e.g., \texttt{"- In step 3..."}). The fallback regex matched the trailing bullet index \texttt{3} instead of the revenue answer. Under structural equivalence, all three rollouts are mathematically correct, establishing 100/100 (100.0\%) accuracy for both checkpoints.

\subsection{Behavioral Trace Decomposition: Continuation Inertia}
We decompose all 200 recovery rollouts across five explicit process dimensions:
\begin{itemize}
    \item \textbf{Final Accuracy ($Y_{\text{acc}}$)}: Exactly 0 / 200 (0.0\%) reached the gold terminal answer.
    \item \textbf{Explicit Error Detection ($Y_{\text{det}}$)}: Exactly 0 / 200 (0.0\%) emitted backtracking cues (*wait, actually, mistake, wrong*).
    \item \textbf{Intermediate Correction ($Y_{\text{corr}}$)}: Exactly 0 / 200 (0.0\%) corrected the corrupted Step 1 product.
    \item \textbf{Faithful Error Propagation ($Y_{\text{prop}}$)}: \textbf{197 / 200 (98.50\%)} adopted the corrupted value $E$ and correctly executed the downstream multiplication $E \times 4$.
    \item \textbf{Attention-Slip Divergence ($Y_{\text{div}}$)}: Exactly 3 / 200 (1.50\%) confused the unit price with the hourly rate.
\end{itemize}

\subsection{Finite-Sample Exact Confidence Bounds}
Because asymptotic inference collapses at zero events ($k=0$), we compute exact Clopper-Pearson binomial confidence bounds \cite{clopper1934use}:
$$p_U = 1 - (\alpha/2)^{1/n}, \quad p_{U, \text{one-sided}} = 1 - \alpha^{1/n}$$
As reported in Table \ref{tab:clopper}, assuming 200 independent trials yields a two-sided 95\% interval of $[0.000\%, 1.828\%]$ (one-sided upper bound: $1.487\%$). However, because rollouts within problem clusters share prefixes, treating the 20 problem clusters as the sampling unit ($C=20$) provides a conservative cluster-level sensitivity bound of $\le 13.911\%$ (two-sided interval: $[0.000\%, 16.843\%]$).

\begin{table}[htbp]
\caption{Clopper-Pearson Exact Confidence Bounds for Zero Observed Recovery ($k=0$)}
\label{tab:clopper}
\centering
\begin{tabular}{lccc}
\toprule
\textbf{Aggregation Tier} & \textbf{Effective $N$} & \textbf{One-Sided 95\% Upper} & \textbf{Two-Sided 95\% CI} \\
\midrule
Rollout Level (Total) & 200 & 1.49\% & $[0.00\%, 1.83\%]$ \\
Model Level (Per Checkpoint) & 100 & 2.95\% & $[0.00\%, 3.62\%]$ \\
\textbf{Cluster Level (Sensitivity)} & \textbf{20} & \textbf{13.91\%} & $\mathbf{[0.00\%, 16.84\%]}$ \\
\bottomrule
\end{tabular}
\end{table}

\section{Discussion: Mechanistic Nature of Continuation Inertia}
Why do small language models exhibit near-total continuation inertia?
Autoregressive language models are trained via next-token prediction to model the joint distribution $P(x_t \mid x_{<t})$. When presented with a prompt prefix containing an arithmetic error, the attention mechanism treats all prefix tokens as grounded context. In standard pretraining and instruction tuning, the model is penalized for contradicting prompt premises. Consequently, the model's most probable trajectory is to execute subsequent derivations that are logically consistent with the preceding tokens, even when those tokens contain an objective arithmetic flaw.

This demonstrates that without explicit verification prompting or external feedback tools (such as CRITIC or Reflexion), autoregressive models function as \textbf{faithful continuation engines} rather than autonomous error detectors.

\section{Limitations}
\begin{itemize}
    \item \textbf{Model Scale}: Evaluated on two 1.5B checkpoints; larger models (7B--72B) or reasoning models with internal test-time search (e.g., DeepSeek-R1) may exhibit different behaviors.
    \item \textbf{Task Domain}: Restricted to 20 elementary arithmetic rate-multiplication problems; does not evaluate open-domain theorem proving or multi-branch logic.
    \item \textbf{Decoding Scaffolding}: Evaluated under unprompted continuation without reflective scaffolding (*"check your work"*) or tool interaction.
    \item \textbf{Observational Models}: Checkpoint comparisons remain observational due to chat template differences.
    \item \textbf{Statistical Power}: With $C=20$ problem clusters, the study has low power to detect subtle recovery rates ($|D| < 0.10$).
    \item \textbf{Human Validation}: Internal blinded human audit of 40 rollouts showed 100\% concordance with structural parsing, but multi-rater crowdsourced adjudication remains pending future work.
\end{itemize}

\section{Reproducibility \& Open Source}
All experimental materials, code, and neural rollout logs are permanently open-sourced at \url{https://github.com/shamddd/recovery_eval}. 
To reproduce all statistical tables in under 1 second:
\begin{verbatim}
git clone https://github.com/shamddd/recovery_eval.git
cd recovery_eval
python3 reproduce_results.py
python3 -m pytest tests/ -v
\end{verbatim}

\section{Conclusion}
We introduced \textbf{RecoveryEval}, an open-source evaluation framework for language model reasoning trajectories. By eliminating control answer leakage, 64-token truncation censoring, and parser artifacts, we demonstrated that the historical negative interaction was an evaluation artifact. Small language models exhibit near-total continuation inertia ($98.5\%$), faithfully propagating injected errors rather than self-correcting. Our findings highlight the necessity of state-matched, provenance-aware infrastructure for reasoning diagnostics.

\section*{Declarations \& AI Assistance Disclosure}
In the preparation of this manuscript, the author utilized automated AI research tools (including Google Antigravity and Large Language Models) to assist with automated test suite refactoring, simulated adversarial peer review critiques, and LaTeX formatting. All primary neural rollouts were generated on independent GPU hardware (RunPod RTX 4090) using official checkpoints, and all statistical counts, exact bounds, and interpretations were verified against raw logs. The human author takes full scientific responsibility for the work. In compliance with publishing standards, no AI system is listed as an author or co-author.

\bibliographystyle{IEEEtran}
\bibliography{references}

\end{document}
